\begin{table}[pos=!htbp,align=\centering]
\caption{Final holdout performance of the CV-selected model at the 15-basis-point depeg threshold. Probability cutoffs are selected separately on each preceding validation fold to maximize event recall within 2–24 hours under a 5-false-alerts-per-month budget. The specification is fixed by CV event recall, with lower false-alert burden breaking ties; no holdout outcome selects the model or cutoff. The selected model detects 3 of 5 eligible holdout events. Timely warnings precede onset by 2–24 hours. Recall and pooled median-lead confidence intervals resample unique realized events; alert-burden intervals resample calendar blocks. Alert burden pools false-alert counts over fully observed follow-up hours. Above-threshold crossings form one episode until 12 consecutive observed below-threshold hours reset it. The first alert requires the same observed quiet history; missing hours or scores break that history. All timely episodes are recognized, while each event contributes once to recall. AUC and AUPRC are secondary row-level ranking metrics. Best displayed values in each metric are highlighted in bold red and second-best distinct displayed values in underlined blue.}
\label{tab:holdout_model_results_15}
\resizebox{\textwidth}{!}{
\begin{tabular}{llrrrrrr}
\toprule
Model & $\alpha$ & Event recall & False alerts/month & Median lead (h) & AUC & AUPRC & Brier score \\
\midrule
\multirow{1}{*}{Kumo} & 1 & \textcolor{red}{\textbf{0.600 [0.200, 1.000]}} & \textcolor{red}{\textbf{2.55 [1.66, 3.55]}} & \textcolor{red}{\textbf{22.0 [14.0, 24.0]}} & \textcolor{red}{\textbf{0.889}} & \textcolor{red}{\textbf{0.066}} & \textcolor{red}{\textbf{0.012}} \\
\bottomrule
\end{tabular}
}
\end{table}
